% Options for packages loaded elsewhere
\PassOptionsToPackage{unicode}{hyperref}
\PassOptionsToPackage{hyphens}{url}
\documentclass[
]{article}
\usepackage{xcolor}
\usepackage[margin=1in]{geometry}
\usepackage{amsmath,amssymb}
\setcounter{secnumdepth}{-\maxdimen} % remove section numbering
\usepackage{iftex}
\ifPDFTeX
  \usepackage[T1]{fontenc}
  \usepackage[utf8]{inputenc}
  \usepackage{textcomp} % provide euro and other symbols
\else % if luatex or xetex
  \usepackage{unicode-math} % this also loads fontspec
  \defaultfontfeatures{Scale=MatchLowercase}
  \defaultfontfeatures[\rmfamily]{Ligatures=TeX,Scale=1}
\fi
\usepackage{lmodern}
\ifPDFTeX\else
  % xetex/luatex font selection
\fi
% Use upquote if available, for straight quotes in verbatim environments
\IfFileExists{upquote.sty}{\usepackage{upquote}}{}
\IfFileExists{microtype.sty}{% use microtype if available
  \usepackage[]{microtype}
  \UseMicrotypeSet[protrusion]{basicmath} % disable protrusion for tt fonts
}{}
\makeatletter
\@ifundefined{KOMAClassName}{% if non-KOMA class
  \IfFileExists{parskip.sty}{%
    \usepackage{parskip}
  }{% else
    \setlength{\parindent}{0pt}
    \setlength{\parskip}{6pt plus 2pt minus 1pt}}
}{% if KOMA class
  \KOMAoptions{parskip=half}}
\makeatother
\usepackage{graphicx}
\makeatletter
\newsavebox\pandoc@box
\newcommand*\pandocbounded[1]{% scales image to fit in text height/width
  \sbox\pandoc@box{#1}%
  \Gscale@div\@tempa{\textheight}{\dimexpr\ht\pandoc@box+\dp\pandoc@box\relax}%
  \Gscale@div\@tempb{\linewidth}{\wd\pandoc@box}%
  \ifdim\@tempb\p@<\@tempa\p@\let\@tempa\@tempb\fi% select the smaller of both
  \ifdim\@tempa\p@<\p@\scalebox{\@tempa}{\usebox\pandoc@box}%
  \else\usebox{\pandoc@box}%
  \fi%
}
% Set default figure placement to htbp
\def\fps@figure{htbp}
\makeatother
\setlength{\emergencystretch}{3em} % prevent overfull lines
\providecommand{\tightlist}{%
  \setlength{\itemsep}{0pt}\setlength{\parskip}{0pt}}
\usepackage{bookmark}
\IfFileExists{xurl.sty}{\usepackage{xurl}}{} % add URL line breaks if available
\urlstyle{same}
\hypersetup{
  pdftitle={Welcome to NRES 470/670},
  pdfauthor={Applied Population Ecology, Spring 2026},
  hidelinks,
  pdfcreator={LaTeX via pandoc}}

\title{Welcome to NRES 470/670}
\author{Applied Population Ecology, Spring 2026}
\date{}

\begin{document}
\maketitle

\textbf{Instructor:} Kevin Shoemaker\\
\textbf{Office:} Fleischmann Agriculture, room 220e\\
\textbf{Phone:} 775-682-7449\\
\textbf{Email}: kevinshoemaker\_at\_unr\_dot\_edu\\
\textbf{Office hours}: Wednesdays from 11am to noon in FA 220e\\
\textbf{Course Website}:
\url{https://kevintshoemaker.github.io/NRES-470/}

\textbf{Teaching Assistant} Jerrod Merrell
(jmerrell\_at\_unr\_dot\_edu)\\
\textbf{TA Office hours}: TBD

\subsubsection{Course Meeting Times}\label{course-meeting-times}

\textbf{Lecture \& Discussion}: M, W at 10am (50 mins) in FA 301
(Fleischmann Ag Building)\\
\textbf{Lab}: F at 1pm (2 hrs 45 mins) in FA 301

\subsubsection{Course description}\label{course-description}

This class explores population ecology theory and how concepts of
population ecology can be used to inform the conservation and management
of natural populations and ecosystems. We emphasize practical approaches
to problem-solving in ecology, conservation, and wildlife management
using simulation models and statistics. Topics include Population
Viability Analysis (PVA), limits to population growth, metapopulation
ecology, species interactions (competition and predation), threats to
wild populations, wildlife management and more. Laboratory exercises
provide students with hands-on experience with wildlife population
models and their practical applications in wildlife ecology and
management.

\subsubsection{Prerequisites}\label{prerequisites}

\begin{itemize}
\tightlist
\item
  BIOL 314 or NRES 217 (Ecology)\\
\item
  NRES 310 (Wildlife Ecology and Management)\\
  NOTE: this course is a prerequisite for NRES 488 (Dynamics and
  Management of Wildlife Populations) and is designed to complement NRES
  421 (Conservation Biology).
\end{itemize}

\subsubsection{Texts}\label{texts}

\begin{itemize}
\tightlist
\item
  \href{https://www.amazon.com/Primer-Ecology-Fourth-Nicholas-Gotelli/dp/0878933182}{Gotelli,
  N. J. A primer of ecology}\\
\item
  Additional readings will be assigned for discussion periodically.
\end{itemize}

\subsubsection{Software}\label{software}

\begin{itemize}
\tightlist
\item
  \href{https://insightmaker.com/}{InsightMaker- web-based systems
  modeling tool}(free, web-based, need account)\\
\item
  \href{https://cran.r-project.org/}{R- software for statistical
  computing and graphics} (free installation)\\
\item
  Spreadsheet software (MS Excel, Google Sheets or equivalent)
\item
  Top Hat (interactive classroom software- invitations should have been
  emailed to you)
\end{itemize}

\subsubsection{Student Learning
outcomes}\label{student-learning-outcomes}

SLO 1. Explain how and why simulation models are used by ecologists and
wildlife professionals.\\
SLO 2. Apply tools such as population viability analysis (PVA) and
metapopulation models to address the conservation and management of wild
populations.\\
SLO 3. Perform basic statistics, data visualization, simulation modeling
and model validation with Excel, the statistical computing language `R',
and the web-based software, InsightMaker.\\
SLO 4. Critically evaluate the strength of inferences drawn from
ecological simulation models using tools such as sensitivity analysis.\\
SLO 5. Explain how species interactions can influence population
dynamics (e.g., predictions of species range shifts).\\
SLO 6. Communicate original research in applied population and community
ecology via professional-style oral and written presentations.

\subsubsection{Grading:}\label{grading}

The course grade will be based on the following components:

Lab exercises (7 total) 20\%\\
Lectures/participation 10\%\\
Group project 30\%\\
Midterm exam \# 1 (date TBD) 10\%\\
Midterm exam \# 2 (date TBD) 10\%\\
Final exam 20\%

Grading scale: A (100 to 93), A- (92 to 90), B+ (89 to 87), B (86 to
83), B- (82 to 80), C+ (79 to 77), C (76 to 73), C- (72 to 70), D+ (69
to 67), D (66 to 63), D- (62 to 60), F (below 60).

\subsubsection{Exams:}\label{exams}

There will be two midterm exams and a final exam, all of which will be
cumulative, covering all course material covered up to the week prior to
the exam. These will consist of multiple-choice, short-answer questions,
and essay questions requiring synthesis of key concepts.

\subsubsection{Lectures}\label{lectures}

Lecture grades will be based primarily on participation and short
in-class quizzes (via Top Hat). Participation is essential to the
learning process (and to our mutual enjoyment of this class). Learning
is not a passive process; students are expected to engage with the
material in class rather than simply listen and take notes. You should
be prepared in class to ask questions, to answer questions, and to
engage in problem-solving activities.

\subsubsection{Labs}\label{labs}

Lab exercises will focus on applying concepts and methods introduced in
lectures, and will involve real-world problems in wildlife conservation
and management wherever possible. Graded lab assignments will involve
figures, tables, InsightMaker models and R code (when applicable) and
responses to questions in short-answer format.

\subsubsection{Final group project}\label{final-group-project}

Students will work in groups of 3-4 to perform a population viability
analysis (PVA) to rank conservation or management actions for a species
of conservation concern (species of your choice!). Grading will be based
on finished products (written and oral presentations) as well as
participation and peer evaluations.

\subsubsection{Graduate credit (for students enrolled in NRES
670)}\label{graduate-credit-for-students-enrolled-in-nres-670}

Graduate students will be subject to additional expectations in order to
receive graduate credit for this course. In particular, graduate
students will be expected to develop an original lecture and lead an
original lab activity related to a topic relevant to wildlife population
ecology. Graduate students will also be expected to achieve a deeper
understanding of the course material, and therefore will be expected to
participate as leaders in discussions and lab activities.

\subsubsection{Make-up policy and late
work:}\label{make-up-policy-and-late-work}

Missed exams and labs cannot be made up, except in the case of
emergencies. If you miss a class meeting, it is your responsibility to
talk to one of your classmates about what you missed. If you miss a lab
meeting, you are still responsible for completing the lab activities and
write-up. Let me or your TA know in advance if you are going to miss
class or lab.

\subsubsection{Top Hat}\label{top-hat}

We will be using the \href{https://tophat.com/}{Top Hat} interactive
learning platform in class. You will be able to submit answers to
in-class questions using Apple or Android smartphones and tablets,
laptops, or through text message.

Top Hat is free of charge for UNR students this year! You should all
have received an invitation via email. You can also enroll using the
join code 866526.

\subsection{University Policies}\label{university-policies}

\subsubsection{Statement on Academic
Dishonesty}\label{statement-on-academic-dishonesty}

The University Academic Standards Policy defines academic dishonesty,
and mandates specific sanctions for violations. See the University
Academic Standards policy:
\href{https://www.unr.edu/administrative-manual/6000-6999-curricula-teaching-research/6502-academic-standards}{UAM
6,502}

\subsubsection{Statement of Disability
Services}\label{statement-of-disability-services}

Any student with a disability needing academic adjustments or
accommodations is requested to speak with the instructor or the
\href{https://www.unr.edu/drc}{Disability Resource Center} (Pennington
Achievement Center Suite 230) as soon as possible to arrange for
appropriate accommodations. This course may leverage 3rd party
web/multimedia content, if you experience any issues accessing this
content, please notify your instructor.

\subsubsection{Statement on Audio and Video
Recording}\label{statement-on-audio-and-video-recording}

Surreptitious or covert video-taping of class or unauthorized audio
recording of class is prohibited by law and by Board of Regents policy.
This class may be videotaped or audio recorded only with the written
permission of the instructor. In order to accommodate students with
disabilities, some students may have been given permission to record
class lectures and discussions. Therefore, students should understand
that their comments during class may be recorded.

Class sessions may be audio-visually recorded for students in the class
to review and for enrolled students who are unable to attend live to
view. Students who participate with their camera on or who use a profile
image are consenting to have their video or image recorded. If you do
not consent to have your profile or video image recorded, keep your
camera off and do not use a profile image. Students who un-mute during
class and participate orally are consenting to have their voices
recorded. If you do not consent to have your voice recorded during
class, keep your mute button activated and only communicate by using the
``chat'' feature, which allows you to type questions and comments live.

\subsubsection{This is a safe space}\label{this-is-a-safe-space}

The University of Nevada, Reno is committed to providing a safe learning
and work environment for all. If you believe you have experienced
discrimination, sexual harassment, sexual assault, domestic/dating
violence, or stalking, whether on or off campus, or need information
related to immigration concerns, please contact the University's Equal
Opportunity \& Title IX office at 775-784-1547. Resources and interim
measures are available to assist you. For more information, please visit
the Equal Opportunity and Title IX page.

\subsubsection{Statement on Campus Closures or
Delays}\label{statement-on-campus-closures-or-delays}

In the event of class cancellations or delays caused by inclement
weather conditions, fire/smoke conditions, or other unforeseen
emergencies, the safety and well-being of students are the University's
top priority. Official notifications will be disseminated through the
University website and other official channels with details related to
any campus delays or closures. In the event of a campus closure, you
will be informed as to whether the class will be offered remotely or if
it will be canceled. If the class is cancelled, you will receive
information on how to address any missed course content. Students facing
significant impacts due to these events are encouraged to communicate
with their instructor for potential accommodations.

\subsubsection{Statement on Failure to Comply with Policy (including as
outlined in this Syllabus) or Directives of a University
Employee:}\label{statement-on-failure-to-comply-with-policy-including-as-outlined-in-this-syllabus-or-directives-of-a-university-employee}

In accordance with section 6,502 of the University Administrative
Manual, a student may receive academic and disciplinary sanctions for
failure to comply with policy, including this syllabus, for failure to
comply with the directions of a University Official, for disruptive
behavior in the classroom, or any other prohibited action. ``Disruptive
behavior'' is defined in part as behavior, including but not limited to
failure to follow course, laboratory or safety rules, or endangering the
health of others. A student may be dropped from class at any time for
misconduct or disruptive behavior in the classroom upon recommendation
of the instructor and with approval of the college dean. A student may
also receive disciplinary sanctions through the Office of Student
Conduct for misconduct or disruptive behavior, including endangering the
health of others, in the classroom. The student shall not receive a
refund for course fees or tuition.

\end{document}
